\section{Cloud-Eigenschaften}
\subsection{Skalierung – Scalability}
Skalierung im Cloud Kontext bedeutet Rechenleistung, Speicher und Netzwerkbandbreite dynamisch an sich ändernde Workloads anzupassen. Es wird in zwei Arten unterteilt.
\subsubsection{Horizontale Skalierung}
Bei horizontaler Skalierung wird ein Workload auf mehreren Servern ausgeführt bzw ausgeweitet.
\subsubsection{Vertikale Skalierung}
Ein bestehender Server wird aufgerüstet durch beispielsweise mehr RAM.
\subsection{Elastizität – Elasticity}
Elastizität im Cloud Kontext bedeutet das eine Anwendung automatisch skaliert werden kann.
\subsection{Hochverfügbarkeit – Availability}
Hochverfügbarkeit bedeutet das eine in der Cloud gehostete Anwendung immer erreichbar ist. Wie viele Ausfälle im Jahr prozentual passieren dürfen ist in einem SLA (Service Level Agreement) festgehalten.
\subsubsection{SLA – Service Level Agreement}
In einem SLA ist Uptime, meist als Nines also z.B. 5 Nines --> 99,999, Reaktion und Lösungszeit, Leistungsparameter, Ausstiegsstrategien und Datenportabilität sowie Sicherheit und Compliance geregelt.
